\documentclass[aspectratio=169]{beamer}

% Persian support
\usepackage{xepersian}
\settextfont[Path={./font/}, Scale=1.3]{IRLotus}
\setlatintextfont[Scale=1]{Times New Roman}

% Theme - simpler for faster presentation
\usetheme{Berlin}
\usecolortheme{seahorse}

% Title information
\title{گزارش پروژه نهایی}
\subtitle{درس مبانی هوش مصنوعی - FIS4041}
\author{نیکی مهدیان \and مبینا یوسفی مقدم}
\date{\today}

% Remove navigation symbols
\setbeamertemplate{navigation symbols}{}

% Custom colors for emphasis
\definecolor{important}{RGB}{200,0,0}
\definecolor{highlight}{RGB}{0,100,200}

\begin{document}

% Title slide
\begin{frame}
\titlepage
\end{frame}

% ============================================
% INTRODUCTION
% ============================================
\begin{frame}{معرفی پروژه}
\begin{block}{چهار بخش اصلی}
\begin{enumerate}
\item \textbf{انتخاب ویژگی} با PSO و GA
\item \textbf{خوشه‌بندی} مشتریان مرکز خرید
\item \textbf{یادگیری Q} - بررسی نظری
\item \textbf{خط لوله ML} برای پیش‌بینی وام
\end{enumerate}
\end{block}
\end{frame}

% ============================================
% QUESTION 1
% ============================================
\section{سوال ۱}

\begin{frame}{انتخاب ویژگی - نتایج کلیدی}
\begin{columns}
\column{0.5\textwidth}
\begin{block}{PSO}
\begin{itemize}
\item همگرایی: تکرار ۷
\item ویژگی: \textcolor{important}{\textbf{Credit\_History}}
\item دقت: \textcolor{highlight}{\textbf{۸۰.۶۷\%}}
\end{itemize}
\end{block}

\begin{block}{GA}
\begin{itemize}
\item همگرایی: نسل ۳
\item ویژگی: \textcolor{important}{\textbf{Credit\_History}}
\item دقت: \textcolor{highlight}{\textbf{۸۰.۶۷\%}}
\end{itemize}
\end{block}

\column{0.5\textwidth}
\begin{block}{نتیجه}
\begin{itemize}
\item هر دو الگوریتم به \textbf{همان جواب} رسیدند
\item کاهش ابعاد: \textcolor{highlight}{\textbf{۹۰.۹\%}}
\item (۱۱ ویژگی $\rightarrow$ ۱ ویژگی)
\end{itemize}
\end{block}

\begin{alertblock}{نکته مهم}
Credit\_History به تنهایی می‌تواند ۸۰\% دقت داشته باشد!
\end{alertblock}
\end{columns}
\end{frame}

% ============================================
% QUESTION 2
% ============================================
\section{سوال ۲}

\begin{frame}{خوشه‌بندی - مقایسه}
\begin{table}[h]
\centering
\large
\begin{tabular}{|c|c|}
\hline
\textbf{الگوریتم} & \textbf{Silhouette Score} \\
\hline
K-Means (K=6) & ۰.۴۲۷۴ \\
\hline
Agglomerative (Ward) & ۰.۴۲۰۱ \\
\hline
\textcolor{important}{\textbf{DBSCAN}} & \textcolor{important}{\textbf{۰.۵۲۹۶}} \\
\hline
\end{tabular}
\end{table}

\begin{block}{DBSCAN - بهترین}
\begin{itemize}
\item پارامترها: eps=0.6, min\_samples=10
\item شناسایی: \textbf{۴ خوشه} + \textbf{۳۳\% نویز}
\item \textcolor{highlight}{\textbf{بهترین عملکرد}} در بین همه
\end{itemize}
\end{block}
\end{frame}

% ============================================
% QUESTION 3
% ============================================
\section{سوال ۳}

\begin{frame}{یادگیری Q - خلاصه}
\begin{block}{معادله}
\begin{equation*}
Q_{t+1} = (1-\alpha)Q_t + \alpha(r_t + \gamma \max Q_{t+1})
\end{equation*}
\end{block}

\begin{columns}
\column{0.5\textwidth}
\begin{block}{پارامترها}
\begin{itemize}
\item $\alpha$: نرخ یادگیری
\item $\gamma$: عامل تنزیل
\item $\epsilon$: اکتشاف
\end{itemize}
\end{block}

\column{0.5\textwidth}
\begin{block}{ویژگی}
\begin{itemize}
\item \textbf{Off-policy}
\item یادگیری سیاست بهینه
\item در حین اکتشاف
\end{itemize}
\end{block}
\end{columns}

\begin{block}{مسائل همگرایی}
\begin{itemize}
\item نرخ یادگیری نامناسب $\rightarrow$ برنامه زمان‌بندی تطبیقی
\item اکتشاف ناکافی $\rightarrow$ کاهش $\epsilon$
\end{itemize}
\end{block}
\end{frame}

% ============================================
% FINAL PROJECT
% ============================================
\section{پروژه نهایی}

\begin{frame}{خط لوله ML - مراحل}
\begin{enumerate}
\item \textbf{پیش‌پردازش}: مدیریت مقادیر گم‌شده، کدگذاری، درمان پرت
\item \textbf{کاهش ابعاد}: PCA (۵ مؤلفه، ۶۹.۷۳\% دقت)
\item \textbf{انتخاب مدل}: Random Forest و Gradient Boosting
\item \textbf{تنظیم ابرپارامترها}: Grid Search و Random Search
\item \textbf{ارزیابی}: معیارهای مختلف
\end{enumerate}
\end{frame}

\begin{frame}{نتایج نهایی}
\begin{table}[h]
\centering
\small
\begin{tabular}{|c|c|c|}
\hline
\textbf{معیار} & \textbf{Random Forest} & \textbf{Gradient Boosting} \\
\hline
Precision & \textcolor{highlight}{\textbf{۷۱.۸۴\%}} & ۶۸.۹۱\% \\
\hline
Recall & ۸۷.۰۶\% & \textcolor{highlight}{\textbf{۹۶.۴۷\%}} \\
\hline
F1-Score & ۷۸.۷۲\% & \textcolor{important}{\textbf{۸۰.۳۹\%}} \\
\hline
\end{tabular}
\end{table}

\begin{block}{تحلیل}
\begin{itemize}
\item \textcolor{important}{\textbf{Gradient Boosting}}: بهترین F1 و Recall
\begin{itemize}
\item شناسایی \textbf{۹۶\%} وام‌های خوب
\end{itemize}
\item \textbf{Random Forest}: بهترین Precision
\begin{itemize}
\item کاهش False Positive
\end{itemize}
\end{itemize}
\end{block}
\end{frame}

\begin{frame}{اهمیت ویژگی‌ها}
\begin{block}{مهم‌ترین ویژگی‌ها (Random Forest)}
\begin{enumerate}
\item \textcolor{important}{\textbf{Credit\_History}} (مهم‌ترین)
\item Applicant\_Income
\item Loan\_Amount
\item Coapplicant\_Income
\end{enumerate}
\end{block}

\begin{alertblock}{تأیید یافته}
Credit\_History در پروژه نهایی نیز به عنوان مهم‌ترین ویژگی تأیید شد!
\end{alertblock}
\end{frame}

% ============================================
% CONCLUSION
% ============================================
\section{نتیجه‌گیری}

\begin{frame}{خلاصه دستاوردها}
\begin{block}{نتایج کلیدی}
\begin{itemize}
\item \textbf{انتخاب ویژگی}: کاهش ۹۰.۹\% با حفظ ۸۰.۶۷\% دقت
\item \textbf{خوشه‌بندی}: DBSCAN با ۰.۵۲۹۶ بهترین عملکرد
\item \textbf{یادگیری Q}: تحلیل نظری و مسائل همگرایی
\item \textbf{خط لوله ML}: F1-Score = ۸۰.۳۹\%
\end{itemize}
\end{block}

\begin{block}{ارزش عملی}
\begin{itemize}
\item سیستم پیش‌بینی وام با دقت بالا
\item تقسیم‌بندی مشتریان برای بازاریابی
\item درک عمیق از الگوریتم‌های AI/ML
\end{itemize}
\end{block}
\end{frame}

\begin{frame}
\begin{center}
\Huge \textbf{سپاسگزاریم}

\vspace{1cm}

\Large سوالات؟
\end{center}
\end{frame}

\end{document}

